\section{Introduction}
\label{sec:geometric-results}

A long alternating collision bridge in a dispersing billiard has a stable
local geometry but an exponentially small physical flux.  The natural
asymptotic problem is therefore relative: after normalization by the exact
hyperbolic flux, does the nonlinear bridge law converge on a contact collar
whose size is independent of the number of flights?  A second question is
inverse.  Which part of the contact geometry is encoded by the limiting law,
and how does that answer depend on which components of the physical record
are retained?  A third question is statistical.  When the table itself is
unknown, what is the local experiment generated by endpoint observations
near the moving support boundary?

This paper gives one chain of answers.  The stationary bridge has a uniform
relative determinant factorization into two half-line actions and two
half-line determinant amplitudes.  The scalar onset probabilities determine
complete symmetrized boundary-energy profiles.  If the signed endpoint
positions are retained, however, their support determines the two
\emph{unsymmetrized} half-line actions.  This yields a block-triangular
inverse for all labelled contact jets, including odd jets, without an
even-contact hypothesis and without supplying the contact curvatures.  On
the statistical side, residual-time integration turns the endpoint density
into a linearly vanishing moving-support law.  Its scalar logarithmic
boundary information has a vector extension with an intrinsic information
matrix.  Applied to a finite-dimensional family of unknown billiard tables,
this gives a local asymptotically normal endpoint experiment; the signed
rigidity theorem implies that finitely many positive endpoint windows have
a nonsingular information matrix on every fixed finite contact-jet model.
The earlier fixed-table comparison between a finite bridge and its boundary
model remains as a different, and sharper, statement about observation
coarsening.

\subsection{The nonlinear relative law}
Fix an oriented closest pair.  Write $j$ for the number of alternating
flights, $g$ for the gap and $\gamma>0$ for the associated hyperbolic
exponent.  In signed transverse endpoint coordinates let $W_j(u,v)$ be the
stationary length and put $E_j=W_j-jg$.  The endpoint Hessian remains
nondegenerate as $j\to\infty$, while the physical mixed derivative
$-W_{j,uv}$ is of order $e^{-j\gamma}$.  Absolute convergence of Hessians
is therefore not enough: the estimate must survive division by the small
physical flux.

Theorem~\ref{thm:v8-main-relative} supplies that relative estimate.  There
are strictly convex half-line actions $S_0,S_1$ and positive determinant
amplitudes $B_0,B_1$, all defined on one fixed endpoint neighborhood, such
that every fixed mixed geometric and offset derivative of the normalized
finite bridge converges exponentially.  For $p=j\bmod2$,
\begin{equation}\label{eq:v19-intro-relative-law}
 \frac{2A\sinh(j\gamma)}{d^2}\Pr(E_{j,e}(d))
 =\mathcal F_p(d)+O_{C^k}(\tau^j),
\end{equation}
where
\begin{equation}\label{eq:v19-intro-limit-law}
 \mathcal F_p(d)=\frac{q_p}{\pi d^2}
 \int(d-S_0(u)-S_p(v))_+B_0(u)B_p(v)\,du\,dv,
 \qquad 0<\tau<1.
\end{equation}
The collar $0\le d\le d_*$ is independent of $j$.  Geometric derivatives
in this statement are taken in the channel-centered variables
$t=jg_e(\xi)+d$, with $d$ fixed; this convention is stated explicitly in
Section~\ref{sec:v18-formal-setup}.

The proof normalizes the exact cofactor identity before passage to the
limit.  The two end perturbations are compared in trace norm, and the
residual-time integral is put in common Morse coordinates.  No factor
proportional to the number of interior collisions is introduced.  The
amplitudes also have an intrinsic dynamical interpretation:
Theorem~\ref{thm:v14-intrinsic-density} identifies their transport under the
physical first-hit maps and Corollary~\ref{cor:v15-determinant-product}
identifies them with normalized scalar linearizing densities of the stable
returns.  Scalar linearization itself is classical; the point here is its
exact equality with the physical half-line Fredholm amplitude and the
normalization used by the billiard flux.

\subsection{Signed endpoints determine general labelled contact jets}
If one integrates the endpoint positions and observes only onset
probabilities, the natural inverse is the symmetrized boundary-energy
profile of Section~\ref{sec:v9-compatibility}.  This is a genuine smooth
function-valued invariant, but it cannot record which inverse branch of a
nonsymmetric contact produced a given energy.  That loss is caused by the
observation map, not by the relative boundary law.

Retaining the signed endpoint coordinates removes that ambiguity.  For an
even number of flights whose first and last contacts have labelled type
$b$, the limiting conditional endpoint density is positive precisely on
\begin{equation}\label{eq:v19-intro-endpoint-support}
             S_b(u)+S_b(v)<d.
\end{equation}
Hence the threshold at the slice $v=0$ is exactly $S_b(u)$.  The two
same-type endpoint laws therefore determine the full signed germs of
$S_0,S_1$, before the energy symmetrization is performed.

Theorem~\ref{thm:v19-signed-rigidity} converts these actions into the
contact geometry.  If $a_b=S_b''(0)$, then the physical onset gives $g$ and
\begin{equation}\label{eq:v19-intro-leading-geometry}
 \gamma=\operatorname{arsinh}(g\sqrt{a_0a_1}),\qquad
 c_b=\cosh\gamma\sqrt{a_b/a_{1-b}},\qquad
 \kappa_b=(c_b-1)/g.
\end{equation}
Thus the leading curvatures are recovered, not supplied.  More importantly,
for arbitrary smooth labelled contact graphs
\[
 \psi_b(y)=\frac{\kappa_b}{2}y^2+
            \sum_{n\ge3}\frac{q_{b,n}}{n!}y^n,
\]
the pair $s_n=(S_0^{(n)}(0),S_1^{(n)}(0))^t$ is affine in the new pair of
jets $q_n=(q_{0,n},q_{1,n})^t$ after the lower jets are fixed, with last-jet
block
\begin{equation}\label{eq:v19-intro-jet-block}
 \begin{pmatrix}
 \coth(n\gamma)&\mathfrak r_0^n\csch(n\gamma)\\
 \mathfrak r_1^n\csch(n\gamma)&\coth(n\gamma)
 \end{pmatrix},
 \qquad \mathfrak r_0\mathfrak r_1=1.
\end{equation}
Its determinant is identically
$\coth^2(n\gamma)-\csch^2(n\gamma)=1$.  This gives a recursive inverse at
every finite order, for odd and even jets alike.  There is no denominator
involving $\kappa_0-\kappa_1$ and no reflection symmetry assumption.  In
the analytic class the two signed endpoint-law germs determine both labelled
contact germs; on connected analytic participating boundaries they determine
the boundary images in their fixed contact frames.  The same finite-order
reconstruction from actual long finite bridges has $O_M(\tau^j)$ error by
the relative action estimate.

The earlier even-contact block of
Theorems~\ref{thm:v12-two-contact} and~\ref{thm:v13-two-flight} remains
useful.  It gives explicit inverses already after residual-time and endpoint
coarsening, and the two-flight theorem gives a direct finite-window physical
design.  It is now a special, more coarsened inverse rather than the maximal
rigidity statement of the paper.

\subsection{Boundary information}
Residual-time integration also changes the singularity of the statistical
experiment.  Before integration, the complete tangent record has a uniform
cap density and its finite-versus-boundary discrepancy is an exclusive
support event of size
$q_j=e^{-j\gamma}$.  After residual time is removed, the endpoint density
vanishes linearly on a moving support boundary.  The ordinary score then has
logarithmically divergent variance.

Theorem~\ref{thm:v17-boundary-information} isolates this mechanism.  If
$f_t=a_t(w_t)_+$ has a compact regular support boundary
$\Sigma=\{w_0=0\}$, then
\begin{equation}\label{eq:v19-intro-scalar-boundary-information}
 H^2(f_t,f_0)=\frac{\mathcal I_\Sigma}{4}t^2\log(1/|t|)+O(t^2),
 \qquad
 \mathcal I_\Sigma=
 \int_\Sigma\frac{a_0(\partial_tw_t|_0)^2}{|\nabla w_0|}\,d\sigma.
\end{equation}
For independently thinned observations, with the failure atom retained,
the finite critical likelihood is Gaussian and has the exact total-variation
profile stated in that theorem.  The supercritical separation follows from
multiplicative Hellinger affinity and does not require a false
support-exclusivity estimate.

For the fixed-table billiard finite-versus-boundary pair, determinant-one
whitening gives
\[
 f_\zeta(x)=\frac2\pi
 (1-e^{-\zeta}x_1^2-e^{\zeta}x_2^2)_+,
 \qquad \mathcal I_\Sigma=1,
\]
for unequal curvatures and either parity.  Since
$\zeta_j=2e^{-j\gamma}+O(e^{-3j\gamma})$, the endpoint scale is
\[
 r_j=e^{-2j\gamma}j\gamma
      =q_j^2\log(1/q_j),
\]
up to the displayed exact normalization.  Theorem~\ref{thm:v17-endpoint-physical}
transfers the limiting testing profile to actual positive offsets.  At that
scale the same raw physical preparations have complete-record distance
converging to one, a nontrivial Gaussian endpoint distance, and
success-indicator distance converging to zero.  This strict hierarchy is a
statement about coarsenings of the same physical record; it is not the
unknown-table experiment described next.

\subsection{An unknown-geometry local experiment}
The vector extension in Theorem~\ref{thm:v19-vector-boundary-lan} treats a
local parameter $\theta\in\mathbb R^r$.  If
$f_\theta=a_\theta(w_\theta)_+$ and
$V=D_\theta w_\theta|_0$, the scalar coefficient becomes the intrinsic
matrix
\begin{equation}\label{eq:v19-intro-information-matrix}
 \mathcal J_\Sigma=
 \int_\Sigma\frac{a_0VV^t}{|\nabla w_0|}\,d\sigma.
\end{equation}
For local alternatives $\theta=\delta_nh$ satisfying
$n p_n\delta_n^2\log(1/\delta_n)\to1$, the thinned product experiment is
locally asymptotically normal:
\begin{equation}\label{eq:v19-intro-lan}
 \log\frac{dP_{n,h}^{\otimes n}}{dP_{n,0}^{\otimes n}}
 =h^t\Delta_n-\frac12h^t\mathcal J_\Sigma h+o_P(1),
 \qquad
 \Delta_n\Rightarrow N_r(0,\mathcal J_\Sigma).
\end{equation}
Thus the limiting experiment is the Gaussian shift with information
$\mathcal J_\Sigma$; when the matrix is nonsingular, its inverse is the
usual local asymptotic covariance lower bound for regular estimators.

For an unknown billiard family at a fixed positive offset, the support
function is
\[
 w_{\eta,d}(u,v)=d-S_{0,\eta}(u)-S_{p,\eta}(v),
\]
so \eqref{eq:v19-intro-information-matrix} becomes an explicit integral of
the geometric support velocity along
$S_0(u)+S_p(v)=d$.  This is an experiment in which the table is the local
parameter; it does not compare two known laws on a fixed table.  On every
fixed finite-dimensional labelled contact-jet model, the two same-type
orientations over a small open interval of positive offsets separate every
nonzero tangent direction by
Theorem~\ref{thm:v19-signed-rigidity}.  Compactness then yields a finite set
of positive endpoint windows whose summed boundary information matrix is
positive definite.  The relative law transfers the LAN experiment to
actual long finite bridges whenever the accumulated finite-to-boundary
error vanishes; in raw preparations the corresponding effective sample
size includes the observed success probability, so failed preparations are
charged.

This gives a statistical counterpart to the geometric rigidity theorem:
the same signed support that determines all finite contact jets also
produces a nonsingular logarithmic information matrix on every fixed finite
jet family.  The result is local and finite-dimensional; the separate
profile-acquisition theorems address the genuinely function-valued inverse
problem under smoothness bounds.

\subsection{Position within nonregular statistics and billiard rigidity}
Parameter-dependent support and nonregular likelihood theory are classical.
Akahira~\cite{Akahira1975a,Akahira1975b} and
Ibragimov--Has'minskii~\cite{IbragimovHasminskii} provide early and general
nonregular asymptotic frameworks.  Smith~\cite{Smith1985} treats densities
which vanish as a power at a moving one-dimensional endpoint and, in the
linearly vanishing case, obtains an efficient normal limit at a nonstandard
rate.  Hirano--Porter~\cite{HiranoPorter} develop limits of experiments and
local asymptotic minimax efficiency for parameter-dependent support.  We do
not claim these general phenomena as new.  The literature audit accompanying
this revision separates those results from the statements proved here.

The statistical novelty claimed here is narrower and geometric: the
regular-hypersurface coefficient and its matrix form, the retained-failure
thinning law, the explicit billiard support velocity and determinant
normalization, and the combination with signed-endpoint rigidity which
proves nonsingularity for finite labelled contact-jet families and transfers
the experiment to actual long bridges.

The geometric information set is likewise different from marked-length or
spectral data.  De Simoi--Kaloshin--Leguil~\cite{DSKL},
Finamore--Leguil~\cite{FinamoreLeguil}, and Zelditch~\cite{Zelditch} obtain
rigidity from substantially different global data under their respective
hypotheses.  Here the observations are near-onset physical collision
records in one labelled channel.  The conclusion of
Theorem~\ref{thm:v19-signed-rigidity} is local contact rigidity from signed
endpoint supports; it makes no assertion about unobserved obstacles or the
lattice.

\subsection{Organization}
Part I develops the stationary action, the relative boundary law, scalar
transport and geometric inversion.  The general signed-endpoint rigidity
theorem follows the earlier coarsened two-contact inverse so that its new
odd-jet block can be compared directly with the historical even block.
Part II proves first the scalar and then the vector boundary-information
theorems, specializes them to the fixed-table endpoint hierarchy and the
unknown-geometry local experiment, and retains the Abel and regularized
profile inverses.  The auxiliary proof chain, historical acquisition
schemes, adaptive experiments, minimax side results and original collision
formulas remain part of the manuscript but are collected behind a single
compendium entry point.  No mathematical module used by the active theorem
chain is deleted in this revision.
